\documentclass{article}

\usepackage{bookmark}
\usepackage{graphicx}
\usepackage{hyperref}
\usepackage{amsmath}
\usepackage{amssymb}
\usepackage{amsthm}
\usepackage{babel}
\usepackage[shortlabels]{enumitem}
\usepackage[a4paper, margin=1in]{geometry}

\newtheorem{definition}{Definition}
\newtheorem{example}{Example}
\newtheorem{remark}{Remark}
\newtheorem{exercise}{Excercise}
\newtheorem{corollary}{Corollary}
\newtheorem{proposition}{Proposition}
\newtheorem{recall}{Recall}


\newenvironment{amatrix}[1]{%
  \left(\begin{array}{@{}*{#1}{c}|c@{}}
}{%
  \end{array}\right)
}

\newcommand{\tr}{\textbf{True}}
\newcommand{\fl}{\textbf{False}}

\title{Advanced Group Theory Various Notes}
\author{William Traub}
\date{March 2026}
\begin{document}
\maketitle
\section{Semi-direct product}
This is motivated by the possible groups with order 18.
\begin{definition}
  Direct Product:
  \begin{itemize}
    \item Internal Direct Product
    \begin{enumerate}
      \item $N \trianglelefteq G, H \trianglelefteq G$
      \item $N \cap H = \{ e\}$
      \item $NH = G$
      \item If G is internal direct product of N, H, then $G \simeq N \times H$ and G is isomorphic to the external direct product of H and H.
      \item There exists the maps $\{e\} \to N \leftrightarrow_{p_1}^{i_1} N \times H$
    \end{enumerate}
    \item Semi direct product
    \begin{itemize}
      \item Internal semi direct
      \begin{enumerate}
        \item $N \trianglelefteq G$, $H \leq G$
        \item $N \cap H = \{e\}$
        \item $NH=G$
        \item Theorem: If G is internal semidirect product N,H then it is isomorphic to external semidirect product $H \simeq N \rtimes H, H \to^{\psi} Aut(N)$.
        \end{enumerate}
        \item External semidirect
        \begin{enumerate}
          \item $(a,b)(a',b') = (a (\psi (b))(a'), bb')$
          \item $N \overset{i_1}{\to} N \times H \underset{i_2}{\overset{p_2}{\leftrightarrow}} H$
        \end{enumerate}
        \item Proof: Suppose D is an internal semidirect product $G \overset{f}{\to} N \rtimes H$. $f(g)^3 = f(nh) = (n, h)$. $f(g_1 g_2) \overset{?}{=} f(g_1)f(g_2)$. $g